\documentclass[10pt,a4paper]{amsart}
\usepackage[margin=19mm]{geometry}
\usepackage{amsmath,amssymb,mathtools}
\usepackage[scr=rsfs,cal=euler]{mathalfa}
\usepackage{xfrac}
\usepackage[unicode,hidelinks,bookmarks=false]{hyperref}
\newcommand{\ie}{i.e.\ }
\newcommand{\cf}{cf.\ }
\newcommand{\resp}{respectively\ }
\newcommand{\Subsection}{\subsection}
\providecommand{\nobreakdash}{\nobreak\hbox{-}}
\setlength{\emergencystretch}{2em}
\multlinegap0pt
\usepackage{bm}
\usepackage{bbm}
\usepackage{dsfont}
\usepackage{xspace}
\usepackage{cancel}
\usepackage{mathtools}
\usepackage{amsthm}
\makeatletter
\@ifundefined{Proposition}{\newtheorem{Proposition}{Proposition}}{}
\makeatother
\title[Forming invariant stochastic differential systems with a giv]{Forming invariant stochastic differential systems with a given first integral}
\author{Konstantin A. Rybakov}
\date{Source: arXiv:2601.04865v2 (CC BY 4.0)}
\begin{document}
\maketitle

\noindent\emph{Source and licence.} Forming invariant stochastic differential systems with a given first integral. Version arXiv:2601.04865v2 (CC BY 4.0), distributed under the Creative Commons Attribution 4.0 International licence, as recorded in the arXiv licence metadata for this exact version and shown on the abstract page https://arxiv.org/abs/2601.04865v2. Full rights notice: licenses/arxiv-2601.04865-CC-BY-4.0.txt.\par\smallskip

\noindent\textit{Editorial scope.} Counted unit: the Proposition whose source label is \texttt{prop1} in the article (source0.tex lines 271--279, source label \texttt{prop1}), delivered together with the article's own complete original proof of it (source0.tex lines 281--347, one \texttt{proof} environment of 67 lines). No mathematics is written, repaired or re-derived: statement and proof are the article's own text line for line. Required context retained uncounted: none. Every definition, display and label of those lines is retained unchanged. Omitted from this package and not redistributed: 1--270, 280--280, 348--1172. Nothing inside a retained excerpt is omitted or altered: every retained body is delivered exactly as the article's lines are written, so no omission receipt of any kind accompanies this package. Citations: the retained excerpts carry no citation command at all, so no bibliography entry of the article is transcribed and no reference is redistributed. Reference adaptation: none was needed and none was made; every reference inside the delivered unit resolves inside it, so no unresolved reference or citation is produced. Printed slips kept literal: no printed word, symbol, hypothesis or number of the counted statement or of its proof was altered. Layout and class: the article's own class is \texttt{amsart} with packages including amsmath, amssymb, amsfonts, dsfont, cite, graphicx, xcolor, hyperref; the wrapper is the lane's pinned \texttt{amsart} editorial wrapper at 10pt on A4, which restates the theorem-like environments the retained text needs and adds the article's own macro definitions that the retained text needs (none), with extra wrapper packages (bm, bbm, dsfont, xspace, cancel, mathtools). The delivered numbering is the unit's own. Assets: no retained span contains a figure environment, a graphics-inclusion command, a PDF-page inclusion, a table or a TikZ picture; no image file is copied into this package, the package declares no asset dependency and no image file is redistributed with it. Licence: version arXiv:2601.04865v2 of this article is distributed under the Creative Commons Attribution 4.0 International licence, as recorded in the arXiv licence metadata for this exact version and shown on the abstract page https://arxiv.org/abs/2601.04865v2. A full rights notice is delivered at licenses/arxiv-2601.04865-CC-BY-4.0.txt. Characters: every retained excerpt is pure ASCII, so the delivered engine, XeLaTeX with its default Latin Modern text font, reports no missing character.
\begin{Proposition}\label{prop1}
If $g_2 \neq 0$ and also $g_3 \neq 0$, \dots, $g_{n-1} \neq 0$ for $n > 3$, then vectors $G,N_1,\dots,N_{n-1}$ are linearly independent, and the determinant of the matrix formed by these vectors is equal to $(-1)^{n-1} |G|^2 \pi_n$, where
\begin{equation}\label{eqDefPi}
  \pi_n = \left\{ \begin{aligned}
    & 1 & & \text{for} ~ n = 2, \\
    & g_2 g_3 \dots g_{n-1} & & \text{for} ~ n > 2.
  \end{aligned} \right.
\end{equation}
\end{Proposition}

\begin{proof}
First, we find the determinant of the $(n \times n)$-matrix $N_G$ whose columns are vectors $G,N_1,\dots,N_{n-1}$. The case $n = 2$ is trivial:
\[
  N_G = \left[ \begin{array}{cc}
    g_1 & -g_2 \\
    g_2 & \phantom{-}g_1
  \end{array} \right], \ \ \ \det N_G = -g_1^2 - g_2^2 = -(g_1^2 + g_2^2) = -|G|^2.
\]

Consider the case $n > 2$:
\[
  N_G = \left[
    \begin{array}{cccccc}
      g_1 & g_2 & 0 & \cdots & 0 & 0  \\
      g_2 & -g_1 & g_3 & \cdots & 0 & 0 \\
      g_3 & 0 & -g_2 & \ddots & \vdots & \vdots \\
      \vdots & \vdots & \vdots & \ddots & g_{n-1} & 0 \\
      g_{n-1} & 0 & 0 & \cdots & -g_{n-2} & g_n \\
      g_n & 0 & 0 & \cdots & 0 & -g_{n-1}
    \end{array}
  \right],
\]
using the Laplace expansion along the last row, i.e.,
\[
  \det N_G = (-1)^{n-1} g_n \left|
    \begin{array}{cccccc}
      g_2 & 0 & \cdots & 0 & 0 \\
      -g_1 & g_3 & \cdots & 0 & 0 \\
      0 & -g_2 & \ddots & \vdots & \vdots \\
      \vdots & \vdots & \ddots & g_{n-1} & 0 \\
      0 & 0 & \cdots & -g_{n-2} & g_n
    \end{array}
  \right| - g_{n-1} \left|
    \begin{array}{cccccc}
      g_1 & g_2 & 0 & \cdots & 0 \\
      g_2 & -g_1 & g_3 & \cdots & 0 \\
      g_3 & 0 & -g_2 & \ddots & \vdots \\
      \vdots & \vdots & \vdots & \ddots & g_{n-1} \\
      g_{n-1} & 0 & 0 & \cdots & -g_{n-2}
    \end{array}
  \right|.
\]

The first determinant on the right-hand side is equal to the product of diagonal elements, and the second one is similar in structure to the original determinant but for size $(n-1) \times (n-1)$. Thus, by denoting $D_n = \det N_G$, we obtain the recurrence formula
\begin{equation}\label{eqDnRecurrent}
  D_n = (-1)^{n-1} g_2 g_3 \dots g_{n-1} g_n^2 - g_{n-1} D_{n-1} = (-1)^{n-1} g_n^2 \pi_n - g_{n-1} D_{n-1}.
\end{equation}

Second, we show that
\[
  D_n = (-1)^{n-1} (g_1^2 + g_2^2 + \ldots + g_n^2) \pi_n = (-1)^{n-1} |G|^2 \pi_n,
\]
using mathematical induction.

This formula is valid for $n = 2$, since $D_2 = -|G|^2$. Further, we assume that
\[
  D_{n-1} = (-1)^{n-2} (g_1^2 + g_2^2 + \ldots + g_{n-1}^2) \pi_{n-1}.
\]
Then, in accordance with the recurrent formula~\eqref{eqDnRecurrent}, we have
\begin{align*}
  D_n & = (-1)^{n-1} g_n^2 \pi_n - g_{n-1} (-1)^{n-2} (g_1^2 + g_2^2 + \ldots + g_{n-1}^2) \pi_{n-1} = \\
  & = (-1)^{n-1} g_n^2 \pi_n + (-1)^{n-1} (g_1^2 + g_2^2 + \ldots + g_{n-1}^2) \pi_n = \\
  & = (-1)^{n-1} (g_1^2 + g_2^2 + \ldots + g_n^2) \pi_n = (-1)^{n-1} |G|^2 \pi_n.
\end{align*}

Therefore, if $g_2 \neq 0$ and also $g_3 \neq 0$, \dots, $g_{n-1} \neq 0$ for $n > 3$, then $D_n = \det N_G \neq 0$ and vectors $G,N_1,\dots,N_{n-1}$ are linearly independent. The proposition has been proven.
\end{proof}

\end{document}
